% ═══════════════════════════════════════════════════════════════
% SPANISCH LEHRERHINWEISE LH-03c: Los colores + Adjektivangleichung
% ═══════════════════════════════════════════════════════════════
% Sequenz 3: Mi casa | Doppelstunde 3c
% Fachfarbe: #c41e3a (Karminrot)
% ═══════════════════════════════════════════════════════════════

\documentclass[11pt, a4paper]{article}

% SW-MODUS TOGGLE
\newif\ifswmodus
\swmodusfalse

% PAKETE
\usepackage[utf8]{inputenc}
\usepackage[T1]{fontenc}
\usepackage[ngerman]{babel}
\usepackage{helvet}
\renewcommand{\familydefault}{\sfdefault}

\usepackage[margin=2cm]{geometry}
\usepackage{enumitem}
\usepackage{tabularx}
\usepackage{booktabs}
\usepackage{longtable}

\usepackage{xcolor}
\usepackage{amssymb}
\usepackage{tcolorbox}
\tcbuselibrary{skins}

\usepackage{fancyhdr}
\usepackage{lastpage}
\usepackage{needspace}

% ═══════════════════════════════════════════════════════════════
% FARBEN
% ═══════════════════════════════════════════════════════════════

\definecolor{spanisch-color}{HTML}{c41e3a}
\definecolor{grau-color}{HTML}{888888}
\definecolor{hellgrau-color}{HTML}{f5f5f5}
\definecolor{basis-color}{HTML}{2a7a4b}
\definecolor{standard-color}{HTML}{2c5aa0}
\definecolor{erweiterung-color}{HTML}{7b1fa2}
\definecolor{loesung-color}{HTML}{c62828}

\definecolor{sw-dark}{HTML}{000000}
\definecolor{sw-gray}{HTML}{666666}
\definecolor{sw-white}{HTML}{ffffff}

\ifswmodus
    \colorlet{spanisch}{sw-dark}
    \colorlet{grau}{sw-gray}
    \colorlet{hellgrau}{sw-white}
    \colorlet{basis}{sw-dark}
    \colorlet{standard}{sw-dark}
    \colorlet{erweiterung}{sw-dark}
    \colorlet{loesung}{sw-dark}
\else
    \colorlet{spanisch}{spanisch-color}
    \colorlet{grau}{grau-color}
    \colorlet{hellgrau}{hellgrau-color}
    \colorlet{basis}{basis-color}
    \colorlet{standard}{standard-color}
    \colorlet{erweiterung}{erweiterung-color}
    \colorlet{loesung}{loesung-color}
\fi

\newcommand{\fachfarbe}{spanisch}

% ═══════════════════════════════════════════════════════════════
% VARIABLEN
% ═══════════════════════════════════════════════════════════════

\newcommand{\thema}{Los colores + Adjektivangleichung}
\newcommand{\sequenz}{03}
\newcommand{\block}{c}
\newcommand{\fach}{Spanisch}
\newcommand{\klasse}{7}
\newcommand{\dauer}{90 Min. (Doppelstunde)}
\newcommand{\lerngruppengroesse}{24}
\newcommand{\datum}{\today}

% ═══════════════════════════════════════════════════════════════
% KOPF- UND FUSSZEILE
% ═══════════════════════════════════════════════════════════════

\pagestyle{fancy}
\fancyhf{}
\renewcommand{\headrulewidth}{0.4pt}
\renewcommand{\footrulewidth}{0.4pt}

\lhead{\fach{} -- Klasse \klasse}
\chead{\textbf{LH-\sequenz\block{} (Lehrerhinweise)}}
\rhead{\datum}

\lfoot{}
\cfoot{}
\rfoot{Seite \thepage\ von \pageref{LastPage}}

% ═══════════════════════════════════════════════════════════════
% BEFEHLE
% ═══════════════════════════════════════════════════════════════

\newcommand{\B}{\textcolor{basis}{\textbf{[B]}}}
\newcommand{\Std}{\textcolor{standard}{\textbf{[S]}}}
\newcommand{\E}{\textcolor{erweiterung}{\textbf{[E]}}}
\newcommand{\loesung}[1]{\textcolor{loesung}{\textbf{#1}}}
\newcommand{\es}[1]{\textcolor{\fachfarbe}{\textbf{#1}}}

\newtcolorbox{kurzplan}{
    colback=hellgrau,
    colframe=\fachfarbe,
    colbacktitle=white,
    coltitle=\fachfarbe,
    boxrule=1pt,
    arc=0pt,
    fonttitle=\bfseries,
    attach boxed title to top left={yshift=-2mm, xshift=5mm},
    boxed title style={boxrule=0pt, colback=white},
    title=Kurzplan
}

\newtcolorbox{differenzierung}{
    colback=white,
    colframe=grau,
    colbacktitle=white,
    coltitle=grau,
    boxrule=0.5pt,
    arc=0pt,
    fonttitle=\bfseries,
    attach boxed title to top left={yshift=-2mm, xshift=5mm},
    boxed title style={boxrule=0pt, colback=white},
    title=Differenzierung
}

\newtcolorbox{fehlerbox}{
    colback=white,
    colframe=loesung,
    colbacktitle=white,
    coltitle=loesung,
    boxrule=0.5pt,
    arc=0pt,
    fonttitle=\bfseries,
    attach boxed title to top left={yshift=-2mm, xshift=5mm},
    boxed title style={boxrule=0pt, colback=white},
    title=Häufige Fehler
}

\newtcolorbox{materialbox}{
    colback=white,
    colframe=\fachfarbe,
    colbacktitle=white,
    coltitle=\fachfarbe,
    boxrule=0.5pt,
    arc=0pt,
    fonttitle=\bfseries,
    attach boxed title to top left={yshift=-2mm, xshift=5mm},
    boxed title style={boxrule=0pt, colback=white},
    title=Material-Checkliste
}

% ═══════════════════════════════════════════════════════════════
% DOKUMENT
% ═══════════════════════════════════════════════════════════════

\begin{document}

% KOPFBEREICH
\begin{center}
    {\Large\bfseries\textcolor{\fachfarbe}{Lehrerhinweise: \thema}}\\[0.3em]
    {\normalsize LH-\sequenz\block{} | \dauer}
\end{center}

\vspace{10pt}

% ═══════════════════════════════════════════════════════════════
% 1. KURZPLAN
% ═══════════════════════════════════════════════════════════════

\section*{1. Stundenverlauf (Kurzplan)}

\textbf{1. Stunde (45 Min.) -- Schwerpunkt: Farben lernen}

\begin{tabularx}{\textwidth}{|c|l|X|l|}
\hline
\textbf{Zeit} & \textbf{Phase} & \textbf{Inhalt} & \textbf{Material} \\
\hline
0--5 & Einstieg & Farbige Gegenstände zeigen: ``¿De qué color es?'' & Objekte \\
\hline
5--15 & Input & 10 Farben einführen, Tafelanschrieb, chorisches Sprechen & Farbkarten \\
\hline
15--25 & Übung 1 & Farben-Drill: Karten zeigen, schnell antworten & Karten \\
\hline
25--35 & Übung 2 & AB-03c Merkkasten 1 ausfüllen & AB \\
\hline
35--45 & Sicherung & Farben an Bildern abfragen, Überleitung Adj.-Angleichung & Bilder \\
\hline
\end{tabularx}

\vspace{1em}

\textbf{2. Stunde (45 Min.) -- Schwerpunkt: Adjektivangleichung}

\begin{tabularx}{\textwidth}{|c|l|X|l|}
\hline
\textbf{Zeit} & \textbf{Phase} & \textbf{Inhalt} & \textbf{Material} \\
\hline
0--5 & Aktivierung & Schnelles Farben-Quiz mit Bildern & Bilder \\
\hline
5--15 & Input & Adjektivangleichung erklären: -o/-a, Tabelle, Merkkasten 2 & Tafel \\
\hline
15--25 & Übung 1 & AB-03c Aufgabe 2: Adjektive in richtiger Form einsetzen & AB \\
\hline
25--35 & Übung 2 & PA: Zimmer beschreiben ``En mi habitación hay...'' & -- \\
\hline
35--40 & Produktion & AB-03c Aufgabe 4: ``Mi casa ideal'' beginnen & AB \\
\hline
40--45 & Abschluss & Zusammenfassung, HA erklären, Ausblick 3d & Tafel \\
\hline
\end{tabularx}

% ═══════════════════════════════════════════════════════════════
% 2. DIFFERENZIERUNG
% ═══════════════════════════════════════════════════════════════

\section*{2. Differenzierung}

\begin{differenzierung}
\textbf{Legende:} \B{} Basis (BR) | \Std{} Standard (MR) | \E{} Erweiterung (GY)

\vspace{0.5em}

\textbf{WICHTIG:} Diese Marker sind NUR für die Lehrkraft. Auf Schülermaterial erscheinen KEINE Niveaumarkierungen!
\end{differenzierung}

\vspace{1em}

\begin{tabularx}{\textwidth}{|l|X|X|X|}
\hline
& \B{} \textbf{Basis} & \Std{} \textbf{Standard} & \E{} \textbf{Erweiterung} \\
\hline
\textbf{Aufgabe 1} & Zuordnung mit Farbkarten als Hilfe & Selbstständige Zuordnung & + eigene Beispiele nennen \\
\hline
\textbf{Aufgabe 2} & Mit Tabelle (Merkkasten 2) & Ohne Tabelle & + Sätze bilden \\
\hline
\textbf{Aufgabe 3} & Vorgegebene Satzanfänge & Mit Struktur-Hilfe & Freie Formulierung \\
\hline
\textbf{Aufgabe 4} & 2 Zimmer, 2 Farben & 3 Zimmer, 4 Farben & 4+ Zimmer, detailliert \\
\hline
\end{tabularx}

\vspace{1em}

\textbf{Konkrete Hinweise:}
\begin{itemize}
    \item \B{} SuS mit LRS: Zeitzugabe (+10 Min.), Farbkarten zur Unterstützung
    \item \B{} DaZ-SuS: Farben vorab auf Deutsch-Spanisch besprechen
    \item \E{} Leistungsstarke: Zusatzaufgabe ``Describe la casa de tus sueños'' (8+ Sätze)
\end{itemize}

% ═══════════════════════════════════════════════════════════════
% 3. HÄUFIGE FEHLER
% ═══════════════════════════════════════════════════════════════

\section*{3. Häufige Fehler}

\begin{fehlerbox}
\begin{tabularx}{\textwidth}{|l|X|X|}
\hline
\textbf{Fehler} & \textbf{Beispiel} & \textbf{Reaktion} \\
\hline
Adjektiv VOR Nomen & *\es{rojo sofá} statt \es{sofá rojo} & Merkspruch: ``Im Spanischen kommt die Farbe NACH dem Ding'' \\
\hline
Fehlende Angleichung & *\es{la mesa rojo} statt \es{la mesa roja} & Tabelle zeigen: ``la = feminin $\rightarrow$ -a'' \\
\hline
Unveränderliche verändern & *\es{verda} statt \es{verde} & Regel: ``Nur -o ändert sich! Verde bleibt verde.'' \\
\hline
marrón falsch & *\es{marrona} statt \es{marrón} & Eselsbrücke: ``Marrón ist stur, bleibt immer marrón'' \\
\hline
Aussprache rr & /maron/ statt /marrón/ & Drill: rr = rollend, Kontrastübung: pero vs. perro \\
\hline
\end{tabularx}
\end{fehlerbox}

% ═══════════════════════════════════════════════════════════════
% 4. MATERIAL-CHECKLISTE
% ═══════════════════════════════════════════════════════════════

\section*{4. Material-Checkliste}

\begin{materialbox}
\begin{itemize}[label=$\square$]
    \item AB-\sequenz\block.pdf (\lerngruppengroesse{} Kopien)
    \item ML-\sequenz\block.pdf (1× für Lehrkraft)
    \item Lehrwerk: Qué pasa 1, S. 88-90
    \item Farbkarten (10 Stück, laminiert: rojo, azul, verde, amarillo, negro, blanco, marrón, gris, rosa, naranja)
    \item Bildkarten: Möbel in verschiedenen Farben (sofá, mesa, silla, cama, armario)
    \item Tafelmarker / Kreide (farbig!)
    \item Optional: LP-\sequenz\block.html (Tablets/PCs)
    \item Optional: Farbige Gegenstände aus dem Klassenraum
\end{itemize}
\end{materialbox}

% ═══════════════════════════════════════════════════════════════
% 5. LÖSUNGEN
% ═══════════════════════════════════════════════════════════════

\section*{5. Lösungen AB-\sequenz\block}

\textbf{Merkkasten 1 (Farben):}

\begin{tabular}{ll@{\hspace{2cm}}ll}
\es{rojo} & $\rightarrow$ \loesung{rot} & \es{blanco} & $\rightarrow$ \loesung{weiß} \\
\es{azul} & $\rightarrow$ \loesung{blau} & \es{marrón} & $\rightarrow$ \loesung{braun} \\
\es{verde} & $\rightarrow$ \loesung{grün} & \es{gris} & $\rightarrow$ \loesung{grau} \\
\es{amarillo} & $\rightarrow$ \loesung{gelb} & \es{rosa} & $\rightarrow$ \loesung{rosa} \\
\es{negro} & $\rightarrow$ \loesung{schwarz} & \es{naranja} & $\rightarrow$ \loesung{orange} \\
\end{tabular}

\vspace{1em}

\textbf{Merkkasten 2 (Adjektivangleichung):}
\begin{itemize}
    \item Adjektive stehen \loesung{nach/hinter} dem Nomen
    \item Tabelle: \loesung{-os} (mask. Pl.), \loesung{-a} (fem. Sg.)
    \item Andere Adjektive ändern nur den \loesung{Plural}
\end{itemize}

\vspace{1em}

\textbf{Aufgabe 1: Zuordnung} (4 Punkte)

\begin{tabular}{ll}
\es{rojo} & $\rightarrow$ \loesung{C (rot)} \\
\es{azul} & $\rightarrow$ \loesung{A (blau)} \\
\es{verde} & $\rightarrow$ \loesung{D (grün)} \\
\es{amarillo} & $\rightarrow$ \loesung{B (gelb)} \\
\end{tabular}

\vspace{1em}

\textbf{Aufgabe 2: Adjektiv anpassen} (4 Punkte)

\begin{tabular}{lll}
a) & La mesa es & \loesung{blanca} (fem. Sg. $\rightarrow$ -a) \\
b) & Los sofás son & \loesung{rojos} (mask. Pl. $\rightarrow$ -os) \\
c) & Las sillas son & \loesung{azules} (Pl., azul = unveränderlich $\rightarrow$ -es) \\
d) & El armario es & \loesung{negro} (mask. Sg. $\rightarrow$ -o) \\
\end{tabular}

\vspace{1em}

\textbf{Aufgabe 3: Zimmer beschreiben} (6 Punkte)

Mögliche Antworten:
\begin{itemize}
    \item \loesung{En mi habitación hay una cama blanca.}
    \item \loesung{El armario es marrón.}
    \item \loesung{Las cortinas son verdes.}
    \item \loesung{La mesa es negra y la silla es azul.}
\end{itemize}

\textit{Bewertung: 2P pro korrekten Satz mit Farbe (max. 6P)}

\vspace{1em}

\textbf{Aufgabe 4: Mi casa ideal} (6 Punkte)

Bewertungskriterien:
\begin{itemize}
    \item Mindestens 3 Zimmer genannt: 2P
    \item Mindestens 4 Farbadjektive korrekt verwendet: 2P
    \item Adjektivangleichung korrekt: 1P
    \item Sprachliche Korrektheit: 1P
\end{itemize}

Beispiellösung:
\begin{quote}
\loesung{Mi casa ideal tiene un salón grande con un sofá rojo y una mesa blanca. La cocina es amarilla con armarios marrones. En mi dormitorio hay una cama azul y las paredes son verdes. El baño es blanco y gris.}
\end{quote}

% ═══════════════════════════════════════════════════════════════
% 6. HAUSAUFGABE
% ═══════════════════════════════════════════════════════════════

\section*{6. Hausaufgabe}

\begin{itemize}
    \item Lehrwerk S. 89, Übung 5 (Farben zuordnen)
    \item Vokabeln lernen: 10 Farben + 5 Möbelstücke
    \item AB-03c Aufgabe 4 fertigstellen (falls nicht fertig)
    \item Optional: LP-03c.html bearbeiten (digital)
\end{itemize}

% ═══════════════════════════════════════════════════════════════
% 7. ZUSATZAUFGABEN (für schnelle SuS)
% ═══════════════════════════════════════════════════════════════

\section*{7. Zusatzaufgaben für schnelle SuS}

\begin{enumerate}
    \item \textbf{Farb-Memory:} 20 Karten (10 Farbkarten + 10 Wortkarten) zuordnen
    \item \textbf{Traumzimmer malen:} Zeichnung + Beschreibung auf Spanisch
    \item \textbf{Partnerdiktat:} Partner beschreibt ein Zimmer, anderer zeichnet
    \item \textbf{Farbgedicht:} Kurzes Gedicht mit 5 Farben schreiben
\end{enumerate}

% ═══════════════════════════════════════════════════════════════
% 8. REFLEXIONSNOTIZEN
% ═══════════════════════════════════════════════════════════════

\section*{8. Reflexionsnotizen (nach Durchführung)}

\begin{tabularx}{\textwidth}{|l|X|}
\hline
\textbf{Was lief gut?} & \\[2em]
\hline
\textbf{Was lief nicht?} & \\[2em]
\hline
\textbf{Zeitmanagement} & \\[2em]
\hline
\textbf{Für nächstes Mal} & \\[2em]
\hline
\end{tabularx}

\end{document}
